\section{Method}
\label{sec:method}

\paragraph{Method: \ours{}.} \ours{} is a latent world model whose output is a \emph{motion} of the feature grids it has
observed. For every future step and every patch, it predicts where in a local window of measured features the patch's
content comes from, whether the patch stays or moves, and a correction for content that motion cannot explain
(\cref{fig:method}; the lower row shows each term of \cref{eq:output} on a held-out DROID window).

\begin{figure}[t]
  \centering
  \includegraphics[width=\linewidth]{architecture.pdf}
  \caption{\textbf{\ours{} architecture.} Top: a predictor turns observed features and actions into one query per patch and step; \ours{} decodes all steps in parallel from observed features. Bottom: on a held-out window, each forecast is kept + moved + corrected observed features.}
  \label{fig:method}
\end{figure}

\subsection{Problem setting}
\label{sec:setting}
A frozen encoder $E$ maps an image to a grid of $N=G\times G$ patch features $\Z_t\in\R^{N\times C}$. Given $H$
observed grids $\Z_{-H+1:0}$, the actions between them $\mathbf{a}_{-H+1:-1}$, and a future action sequence
$\mathbf{a}_{0:K-1}$, a world model predicts $\hat\Z_{1:K}$. The forecast $\hat\Z_k$ may depend only on
$\mathbf{a}_{<k}$ (causality). Features are standardised per channel with train-split statistics that are shared
across positions.

\subsection{Architecture}
\label{sec:arch}
\textbf{Memory encoder.} Each observed patch feature is projected to width $d$ and receives a spatial embedding, a frame
embedding and an embedding of the action applied after that frame. $L_e$ transformer layers over all $H\cdot N$ tokens
give memory states $\mathbf{M}\in\R^{H\times N\times d}$.
\textbf{Action-prefix queries.} A unidirectional GRU over embedded actions gives a causal prefix state, to which we add a
horizon embedding to form $\mathbf{c}_k$. Query tokens start from the memory states of the last frame, and $L_d$ decoder
layers apply self-attention across patches and cross-attention to $\mathbf{M}$, each modulated by $\mathbf{c}_k$ through
adaptive layer normalisation with zero-initialised gains \citep{peebles2023dit}. The result is $\h_{k,i}\in\R^d$ for
patch $i$ at step $k$; all $K$ steps are decoded in parallel.

\subsection{Gated transport}
\label{sec:transport}
For query patch $i$, let $\mathcal{W}(i)$ be the $w\times w$ spatial window around $i$ in each of the last $S\le H$
observed frames. Transport weights are
\begin{equation}
  \pi_{k,i}(j) \;=\; \operatorname{softmax}_{j\in\mathcal{W}(i)}\Big(
  \tfrac{1}{\sqrt{d_k}}\,(\mathbf{W}_q\h_{k,i})^{\!\top}(\mathbf{W}_k\mathbf{M}_{j}) + \beta\,\mathbf{1}[j=(0,i)]\Big),
  \label{eq:transport}
\end{equation}
where $\beta$ is a learned identity bias toward the same location in the last frame. Transport moves \emph{measured}
features, not latent states: $\mathbf{T}_{k,i}=\sum_{j}\pi_{k,i}(j)\,\Z_j$. A per-patch gate
$g_{k,i}=\sigma(\mathbf{w}_g^\top\h_{k,i}+b_g)$ and a correction $\mathbf{r}_{k,i}=\mathbf{s}\odot\mathbf{W}_o\h_{k,i}$ give
\begin{equation}
  \hat\Z_{k,i} \;=\; (1-g_{k,i})\,\Z_{0,i} \;+\; g_{k,i}\,\mathbf{T}_{k,i} \;+\; \mathbf{r}_{k,i}.
  \label{eq:output}
\end{equation}
$\mathbf{W}_o$ is zero-initialised and $b_g=-3$, so an untrained model starts close to persistence and learns
departures from it. $\mathbf{s}\in\R^C$ is a learned per-channel scale (\cref{sec:ablations} compares it with a bounded
$\tanh$ correction).

\subsection{Training}
\label{sec:training}
We minimise the mean squared error over all horizons, $\mathcal{L}_{\text{pred}}=\frac1K\sum_k\|\hat\Z_k-\Z_k\|^2$.
Anchoring every forecast to observed features makes it easy to ignore actions: a forecast that stays close to $\Z_0$
scores well on average. For planning we therefore add an action-contrastive hinge that compares the error under the
true future actions with the error under actions borrowed from another sample in the batch:
\begin{equation}
  \mathcal{L}_{\text{act}} = \max\!\big(0,\; m - (\|\hat\Z(\tilde{\mathbf a})-\Z\|^2 - \operatorname{sg}\|\hat\Z(\mathbf a)-\Z\|^2)\big),
  \qquad \mathcal{L}=\mathcal{L}_{\text{pred}}+\lambda\,\mathcal{L}_{\text{act}} .
  \label{eq:contrastive}
\end{equation}
All forecasting and plug-in results use $\lambda{=}0$, like every baseline; $\mathcal{L}_{\text{act}}$ appears only in a planning variant and its ablation.

\subsection{Motion model}
\label{sec:theory}
\begin{wrapfigure}[19]{r}{0.47\linewidth}
  \vspace{-1.2em}
  \centering
  \includegraphics[width=\linewidth]{featurespace.pdf}
  \caption{\textbf{One patch in feature space} under the three forecast routes \ding{172}--\ding{174}.}
  \label{fig:featurespace}
  \vspace{-0.6em}
\end{wrapfigure}
We motivate \cref{eq:output} with a simple generative model. \Cref{fig:featurespace} shows its geometry for one patch:
\ours{} mixes observed vectors (\ding{172}, Prop.~\ref{prop:bayes}), an additive head synthesises the displacement
(\ding{173}, Prop.~\ref{prop:residual}), and recursion feeds errors forward (\ding{174}, Prop.~\ref{prop:compound}).
Suppose the grid at step $k$ arises from motion plus innovation,
\begin{equation}
  \Z_k \;=\; \mathbf{P}_k\,\Z_0 + \bm{\varepsilon}_k,
  \label{eq:motion-model}
\end{equation}
where $\mathbf{P}_k\in\{0,1\}^{N\times N}$ is a row-stochastic \emph{motion} matrix (each patch copies one source
patch $\pi(i)$ at integer-patch offsets; $\mathbf{P}_k=\mathbf{I}$ is a static scene; sub-patch motion is covered by soft weights between neighbouring sources) and $\bm{\varepsilon}_k$ is new content with
$\E[\bm{\varepsilon}_k\mid\mathbf{P}_k,\Z_0,\mathbf{a}]=\mathbf{0}$. Motion is uncertain given the context $c=(\Z_{\le 0},\mathbf{a}_{<k})$.

\begin{proposition}[The optimal forecast is a transport]\label{prop:bayes}
Under \cref{eq:motion-model}, the MSE-optimal forecast is
$\E[\Z_k\mid c]=\bar{\mathbf{P}}_k\Z_0$ with $\bar{\mathbf{P}}_k=\E[\mathbf{P}_k\mid c]$, a row-stochastic matrix.
If the support of each row of $\mathbf{P}_k$ lies in the transport window $\mathcal{W}$, then $\bar{\mathbf{P}}_k\Z_0$
can be approximated to arbitrary precision by the transport of \cref{eq:transport} with $g{=}1$ and
$\mathbf{r}{=}\mathbf{0}$, given logits that the decoder can express.
\end{proposition}

Any sufficiently expressive head could reach this optimum; transport parameterises it exactly, an inductive bias we
test empirically. Real features also change through viewpoint, lighting and deformation, which the correction
$\mathbf{r}$ absorbs. An additive head $\Z_0+\mathbf{W}_o\h$ instead has to write out
the displacement in feature space:

\begin{proposition}[Size of the regression target]\label{prop:residual}
For a deterministic motion $\mathbf{P}$ with sources $\pi(i)$, the residual an additive head must output is
$\mathbf{R}=(\mathbf{P}-\mathbf{I})\Z_0$ with
$\|\mathbf{R}\|_F^2=\sum_i\|\mathbf{z}_{0,\pi(i)}-\mathbf{z}_{0,i}\|^2$, the feature contrast along the motion. The
residual required by \ours{} is $\mathbf{0}$ for every such $\mathbf{P}$ within the window.
\end{proposition}

This concerns what each head must represent, not learnability: an additive head can learn to copy, but its target
grows with motion and feature contrast. On static patches ($\mathbf{P}=\mathbf{I}$) the
identity path is exact, whereas a regression head must output exactly zero, and any error corrupts the static scene;
\cref{tab:regions} shows both effects. Finally, \ours{}, like the Direct baseline, decodes every horizon from measured
features, which avoids the compounding of recursive decoding (proofs in \cref{app:proofs}):

\begin{proposition}[Recursive versus direct decoding]\label{prop:compound}
Let a recursive predictor $f$ be $L$-Lipschitz in its input with one-step error at most $\delta$. Its $k$-step
rollout error is bounded by $\delta\sum_{j=0}^{k-1}L^j$, which grows geometrically for $L>1$. A direct predictor with
per-horizon error $\delta_k$ incurs exactly $\delta_k$, independent of its errors at other horizons.
\end{proposition}

